\section[Harmonic Pullbacks and Transfer]{Harmonic Pullbacks and Subcalibration Transfer}
\label{sec:quartic-application}

Riedler proves regular-locus minimality under umbilic harmonic morphisms
and area minimization for Clifford--Simons quartics
\cite[Theorems~5.1.4 and~5.2.5]{RiedlerThesis25}, and asks about general
minimal cones \cite[Question~5.3.3(2), p.~88]{RiedlerThesis25}.
\Cref{thm:clifford-area-transfer,thm:smallest-clifford-pullback}
control the dilation term with a source-radius factor and determine the
smallest algebra, yielding \cref{cor:iterated-clifford-transfer}.
Quadratic and cubic applications appear in
\cref{sec:two-quadratic,sec:nonradial-examples}.

\subsection{The pullback identity and the dilation term}

For a degree-\(d\) homogeneous polynomial harmonic morphism
\(\phi:\R^N\to\R^m\), \(d\ge1\), the Fuglede--Ishihara
characterization, as recalled in
\cite[Theorem~2.1.2]{RiedlerThesis25},
gives, on its regular set, a dilation \(\lambda>0\) with
\[
 \Delta\phi_i=0,\qquad
 \Gamma(\phi_i,\phi_j)=\lambda^2\delta_{ij}.
\]
In particular, \(\|d\phi\|_{\mathrm{HS}}^2=m\lambda^2\).
Subscripts \(x\) and \(y\) distinguish source and target operators.

\begin{proposition}[Pullback criterion]
\label{prop:pullback-subcalibration}
Suppose that, on the regular set,
\begin{equation}\label{eq:pullback-dilation-condition}
 d\phi\bigl(\nabla_x\lambda^2\bigr)=\eta(x)\phi(x).
\end{equation}
Equivalently,
\(d\phi(\nabla_x\|d\phi\|_{\mathrm{HS}}^2)=m\eta\phi\);
this records the factor \(m\) in the differential-norm formulation
of \cite[Definition~5.1.1(1)]{RiedlerThesis25}.
Let \(H\) be a nonzero homogeneous polynomial of degree \(\ell>0\)
on \(\R^m\), put \(G=H\circ\phi\), and write
\(W_H=|\nabla_yH|^2\).  Then
\begin{align}
 |\nabla_xG|^2&=\lambda^2(W_H\circ\phi),\qquad
 \Delta_xG=\lambda^2(\Delta_yH)\circ\phi,
 \label{eq:pullback-gradient-laplacian}\\
 \cL_x(G)&=\lambda^4\cL_y(H)\circ\phi
       -\frac{\ell\eta}{2}(H W_H)\circ\phi.
 \label{eq:pullback-defect}
\end{align}
Suppose that a constant \(\kappa\in\R\) satisfies
\begin{equation}\label{eq:pullback-kappa-bound}
 \frac{|\phi(x)|^2\eta(x)}{2\lambda(x)^4}\leq\kappa
\end{equation}
at regular points with \(\phi(x)\ne0\), and that
\begin{equation}\label{eq:pullback-target-sign}
 |y|^2H(y)\cL_y(H)(y)
 \geq\kappa\ell H(y)^2W_H(y)
 \qquad(y\ne0).
\end{equation}
Then \(G\cL_x(G)\geq0\) at regular points with \(\phi(x)\ne0\).

If \(G\not\equiv0\) and a closed algebraic set \(Z\subset\R^N\)
of dimension at most \(N-2\) contains the critical set of \(G\),
the singular zeros of \(G\), and the points excluded from this
calculation, then \(\{G<0\}\) is locally perimeter minimizing and
\(\partial\llbracket\{G<0\}\rrbracket\) is locally mass minimizing.
\end{proposition}

\begin{proof}
The chain rule gives \eqref{eq:pullback-gradient-laplacian}.  Moreover,
\begin{align*}
 \tfrac12\Gamma_x(G,|\nabla_xG|^2)
 &=\frac{W_H\circ\phi}{2}
   \left\langle d\phi(\nabla_x\lambda^2),
                   (\nabla_yH)\circ\phi\right\rangle\\
 &\quad+\lambda^4\bigl(\Hess_yH(\nabla_yH,\nabla_yH)\bigr)\circ\phi.
\end{align*}
Euler's identity \(\langle y,\nabla_yH\rangle=\ell H\)
and \eqref{eq:pullback-dilation-condition} prove
\eqref{eq:pullback-defect}.  At a regular point with
\(y=\phi(x)\ne0\), the two assumed inequalities give
\[
 G\cL_x(G)
 \geq\ell H(y)^2W_H(y)
       \left(\frac{\kappa\lambda^4}{|y|^2}
                    -\frac{\eta}{2}\right)\geq0.
\]
The global conclusion follows from
\cref{rem:polynomial-finite-perimeter,cor:gradient-subcalibration},
since the algebraic dimension bound makes \(Z\) locally
\(\cH^{N-1}\)-null.
\end{proof}

Both terms in \eqref{eq:pullback-defect} vanish on a regular minimal
level cone \(H=0\); away from the cone, the dilation term can change
the sign. For example, if \(H\) is linear and
\(\eta>0\), then \(\cL_y(H)=0\) but
\(G\cL_x(G)=-\eta H^2W_H/2<0\) wherever \(H\ne0\).

\begin{remark}[Weighted interpretation]
Condition \eqref{eq:pullback-target-sign} gives the phase inequality
for the normalized gradient with density \(|y|^{-\kappa}\).
Indeed, for \(X_H=\nabla H/|\nabla H|\),
\[
 |y|^{\kappa}\diver_y\bigl(|y|^{-\kappa}X_H\bigr)
 =\frac{\cL_y(H)}{W_H^{3/2}}
       -\frac{\kappa\ell H}{|y|^2W_H^{1/2}}.
\]
It is the density of the Euler--Dierkes--Huisken functional
\(\mathcal E_\alpha(\Sigma)=\int_\Sigma |y|^\alpha\,d\mathcal H^{m-1}\)
with \(\alpha=-\kappa\).
Minimal cones are stationary on their regular parts for area functionals
with smooth positive radial densities. Subcalibration methods also apply
to proving minimality for these functionals; see, for example,
Dierkes~\cite{Dierkes26Euler}.
For the Clifford maps below, \(\alpha=-1/2\); the dilation bound
converts the weighted inequality into ordinary-area comparison.
\end{remark}

For the subclass \(\lambda^2=c|x|^{2d-2}\), with \(c>0\),
Euler's identity gives
\[
 \eta=2d(d-1)c|x|^{2d-4},\qquad
 d^2|\phi(x)|^2=|d\phi_x(x)|^2\leq\lambda^2|x|^2.
\]
Hence \eqref{eq:pullback-kappa-bound} holds with
\(\kappa=(d-1)/d\).

\subsection{Area-minimizing Clifford pullbacks}

Let \(A_1,\ldots,A_m\) be a symmetric Clifford system on \(\R^N\),
where \(m\geq2\):
\[
 A_iA_j+A_jA_i=2\delta_{ij}\Id.
\]
Set \(P_i(x)=\langle A_ix,x\rangle\), \(P=(P_1,\ldots,P_m)\),
\(R=|x|^2\), and \(Q=|P|^2\).  The Clifford relations imply
\begin{equation}\label{eq:Clifford-pullback-identities}
 \Delta P_i=0,\qquad
 \Gamma(P_i,P_j)=4R\delta_{ij},\qquad
 \Gamma(R,P_i)=4P_i.
\end{equation}
For example, a second Clifford matrix interchanges the two eigenspaces
of \(A_i\), so \(\tr A_i=0\). The classification of quadratic
harmonic morphisms by Ou and Wood, including its Clifford-system
and direct-sum description, is recalled in
\cite[Theorem~4.5.3 and Remark~4.5.5]{RiedlerThesis25}.
The maps above are normalized umbilic maps: all positive eigenvalues
of their component matrices equal one; see
\cite[Remark~5.1.12]{RiedlerThesis25}.

\begin{lemma}[Laplacian algebra under Clifford pullback]
\label{lem:Clifford-algebra-pullback}
If \(\cA\subset\R[y_1,\ldots,y_m]\) is a graded Laplacian algebra,
then \(\R[R,P^*\cA]\) is a Laplacian algebra on \(\R^N\).
Moreover, \(P\) is onto, \(Q\leq R^2\), and
\(P^{-1}(0)\setminus\{0\}\), when nonempty, is a smooth
submanifold of codimension \(m\).
For this map, \cref{prop:pullback-subcalibration} applies with
\(\lambda^2=4R\), \(\eta=16\), and \(\kappa=1/2\).
\end{lemma}

\begin{proof}
For target polynomials \(u,v\), the chain rule gives
\[
 \Delta_x(u\circ P)=4R(\Delta_yu)\circ P,\qquad
 \Gamma_x(u\circ P,v\circ P)=4R\Gamma_y(u,v)\circ P.
\]
If \(u\) has degree \(j\), then
\(\Gamma_x(R,u\circ P)=4j(u\circ P)\).
Together with \(\Delta R=2N\) and \(\Gamma(R,R)=4R\), these
identities and the product rule prove the asserted closure.

For \(x\ne0\), the vectors \(A_ix\) are orthogonal and have
length \(|x|\).  Projection of \(x\) onto their span gives
\(Q/R\leq R\).  To prove surjectivity, take \(u\in\Sph^{m-1}\)
and put \(A(u)=\sum u_iA_i\).  This symmetric involution has trace
zero, since every \(A_i\) does.  It therefore has a unit eigenvector
\(x\) with eigenvalue one.  For \(v\perp u\),
the matrices \(A(v)\) and \(A(u)\) anticommute, so
\(\langle A(v)x,x\rangle=0\).  Hence \(P(x)=u\), and scaling
proves that \(P\) is onto.  Finally, \eqref{eq:Clifford-pullback-identities}
shows that \(dP\) has rank \(m\) at every nonzero point.  The
submersion theorem proves the statement about \(P^{-1}(0)\).
Since \(dP(\nabla(4R))=16P\), the last assertion follows from
\(Q/(2R^2)\leq1/2\).
\end{proof}

In particular, if a target homogeneous polynomial \(f\) has degree
at least two and \(\nabla f(y)\ne0\) for \(y\ne0\), then
\[
 |\nabla(f\circ P)|^2=4R(|\nabla f|^2\circ P).
\]
Thus its critical set is \(Z=P^{-1}(0)\), and its zero set is
regular off \(Z\).  The set \(Z\) is algebraic and has dimension
at most \(N-m\), apart from the possible isolated origin.
Consequently it has zero \(\cH^{N-1}\)-measure. Surjectivity preserves
nonemptiness of both phases. These observations apply to the quadratic
and Cartan targets below.

For a nonflat target hypersurface cone with both phases nonempty
and the isolated critical point just assumed, this is also the
geometric singular set:
\[
 \Sing\{f\circ P=0\}=P^{-1}(0).
\]
At every nonzero point of \(P^{-1}(0)\), submersion coordinates
identify the zero set with a product of the target cone and a smooth
factor. The vertex is singular as well: a smooth cone would be a
hyperplane, which cannot be the zero set of an even homogeneous
polynomial taking both signs.

\label{sec:clifford-area-transfer}

A source-radius factor compensates for the dilation term in
\eqref{eq:pullback-defect} in sufficiently large admissible dimensions.

\begin{theorem}[Area minimization under Clifford pullback]
\label{thm:clifford-area-transfer}
Let \(f\) be a nonzero homogeneous polynomial on \(\R^m\), and let
\(Z\subset\{f=0\}\) be a closed conical real algebraic set containing
the origin and having codimension at least two.  Suppose that
\(H\in C^3(\R^m\setminus Z)\cap C^0(\R^m)\) is homogeneous of
degree \(\ell>0\), with
\[
 \operatorname{sgn}H=\operatorname{sgn}f,\qquad
 |\nabla_yH|>0\quad\text{on }\R^m\setminus Z.
\]
Put \(W_H=|\nabla_yH|^2\), and suppose that, off \(Z\),
\[
 H\cL_y(H)\geq0,\qquad H\Delta_yH\geq-MW_H
\]
for some finite \(M\geq0\).  Let \(P:\R^N\to\R^m\) be the
Clifford map above.  If \(a>0\) and
\begin{equation}\label{eq:clifford-transfer-bound}
 N\geq1+4\ell(1+M)+2a(1+M)+\frac{\ell}{a},
\end{equation}
then
\[
 G=R^aH(P),\qquad E=\{f(P)<0\},
\]
defines a two-sided subcalibration: off \(P^{-1}(Z)\),
\(\nabla_xG\ne0\) and \(G\cL_x(G)\geq0\), with strict inequality
in the two open phases.  Consequently \(E\) is locally perimeter
minimizing and \(\partial\llbracket E\rrbracket\) is locally mass
minimizing. If \(N\geq3\) and both phases are nonempty, this current
is strictly area minimizing. If both signs of \(f\) occur in every
neighborhood of each target zero, then
\[
 \operatorname{spt}\partial\llbracket E\rrbracket=\{f(P)=0\}.
\]
In particular, this last condition holds when \(f\) is harmonic.
\end{theorem}

\begin{proof}
Set \(h=H(P)\) and \(b=a(a+2\ell)\).  The chain rule and Euler's
identity give, on \(\R^N\setminus P^{-1}(Z)\),
\begin{align}
 |\nabla_xG|^2
 &=4R^{2a+1}W_H(P)+4bR^{2a-1}h^2,
 \label{eq:clifford-transfer-gradient}\\
 \cL_x(G)
 &=16R^{3a+2}(\cL_yH)(P)
   +8\bigl[a(N-1)-\ell-2b\bigr]R^{3a}hW_H(P)\notag\\
 &\quad+16bR^{3a}h^2(\Delta_yH)(P)
   +8b\bigl[a(N-1)+\ell\bigr]R^{3a-2}h^3.
 \label{eq:clifford-transfer-defect}
\end{align}
For the second formula, use the radial identity
\eqref{eq:radial-master-L} with the source function \(h\), of degree
\(2\ell\), and substitute
\[
\begin{gathered}
 |\nabla_xh|^2=4RW_H(P),\qquad
 \Delta_xh=4R(\Delta_yH)(P),\\
 \cL_x(h)=16R^2(\cL_yH)(P)-8\ell hW_H(P).
\end{gathered}
\]
These identities follow from \eqref{eq:Clifford-pullback-identities}.

Multiplication of \eqref{eq:clifford-transfer-defect} by \(h\) yields
\begin{align*}
 h\cL_x(G)\geq{}&
 8\bigl[a(N-1)-\ell-2b(1+M)\bigr]R^{3a}h^2W_H(P)\\
 &+8b\bigl[a(N-1)+\ell\bigr]R^{3a-2}h^4.
\end{align*}
The first coefficient is nonnegative by
\eqref{eq:clifford-transfer-bound}, and the last term is positive
when \(h\ne0\).  Since \(R^a>0\), this proves the claimed phase
inequality.  Formula \eqref{eq:clifford-transfer-gradient} gives
\(\nabla_xG\ne0\) off \(Z'=P^{-1}(Z)\).

The extension of \(G\) is continuous with the signs of \(f(P)\);
at the origin, use \(H(P(x))=O(|x|^{2\ell})\).
Surjectivity gives \(f(P)\not\equiv0\), hence local finite perimeter.
Since \(P\) is a submersion off the origin, stratifying \(Z\) shows
that \(Z'=P^{-1}(Z)\) has codimension at least two. Thus
\cref{cor:gradient-subcalibration} applies. For \(N\ge3\) and two
nonempty phases, the strict divergence signs give strict minimality
by \cref{thm:strict-minimality-flux}.

The additional phase assumption transfers by submersion away from
the origin and by surjectivity and homogeneity at the origin.
The support conclusion follows from \cref{rem:polynomial-finite-perimeter}.
\end{proof}

For harmonic nonconstant \(f\), \eqref{eq:Clifford-pullback-identities}
and surjectivity make \(F=f\circ P\) nonzero and harmonic.
Area minimization then implies real irreducibility by
\cref{prop:minimizing-irreducibility}. Nonharmonic targets still
require the stated phase hypothesis for full support.

The ratio bound is preserved under iteration:
\[
 \Delta_xG=4R^{a+1}(\Delta_yH)(P)
       +2a(N+2a+4\ell-2)R^{a-1}h,
\]
and \eqref{eq:clifford-transfer-gradient} gives
\begin{align}
 G\Delta_xG+M|\nabla_xG|^2
 ={}&4R^{2a+1}\bigl(H\Delta_yH+MW_H\bigr)(P)\notag\\
 &+\bigl[2a(N+2a+4\ell-2)+4Mb\bigr]R^{2a-1}h^2
 \geq0.
 \label{eq:clifford-transfer-ratio}
\end{align}
Thus the same \(M\) is available for iteration; see
\cref{cor:iterated-clifford-transfer}. In particular, \(a=1\)
preserves polynomial potentials and gives the sufficient bound
\(N\geq5\ell+3+(4\ell+2)M\).

The harmonic radial criterion supplies further target potentials.
Let \(f\) be homogeneous of degree \(e\) on \(\R^m\) and satisfy
the hypotheses of \cref{cor:harmonic-radial-existence}. Set
\(\sigma=(m-1-2e)/2>0\) and \(H=|y|^\sigma f\). Harmonicity gives
\[
 H\Delta_yH
 =\sigma(\sigma+2e+m-2)|y|^{2\sigma-2}f^2\geq0
 \qquad(y\ne0).
\]
Together with the corollary's phase sign and nonvanishing gradient,
this verifies all target hypotheses with \(Z=\{0\}\),
\(\ell=e+\sigma\), and \(M=0\). Thus
\cref{thm:clifford-area-transfer} applies whenever the source-dimension
bound \eqref{eq:clifford-transfer-bound} holds.


\subsection{The smallest algebra and iteration}
\label{subsec:clifford-smallest-algebra}

Write $q(y)=|y|^2$ on the target.

\begin{theorem}[The smallest algebra under Clifford pullback]
\label{thm:smallest-clifford-pullback}
Let $P:\R^N\to\R^m$ be the Clifford map above and let $f$ be a
nonzero essential homogeneous polynomial of degree $s\ge2$ on
$\R^m$, so $V_{\mathrm{cyl}}(f)=\{0\}$. Then
\begin{equation}\label{eq:smallest-clifford-pullback}
 \cA_R(f\circ P)=\R[R,P^*\cA_q(f)].
\end{equation}
If $Q<R^2$ somewhere, this is isomorphic to the polynomial extension
$\cA_q(f)[t]$, with target degrees doubled and $\deg t=2$.
In particular, its minimal number of homogeneous generators is one
more than that of $\cA_q(f)$.
\end{theorem}

\begin{proof}
Set $\cA=\cA_q(f)$ and $\cB=\cA_R(f\circ P)$.
By \cref{lem:Clifford-algebra-pullback}, $\R[R,P^*\cA]$ is Laplacian
and contains $f\circ P$, so $\cB\subseteq\R[R,P^*\cA]$.
For the reverse inclusion put \(\cC=\{u\in\R[y_1,\ldots,y_m]:u\circ P\in\cB\}\).
This is a unital algebra containing $f$. If $u\in\cC$, the chain
rule gives $4R(\Delta u)\circ P=\Delta(u\circ P)\in\cB$.
The quotient is a polynomial, so
\cref{prop:MR-finite-and-division} gives $(\Delta u)\circ P\in\cB$.
Thus $\cC$ is closed under $\Delta$ and, by the product rule, $\Gamma$.

It remains to recover $q$. The recursion
\eqref{eq:MR-higher-contraction} uses only multiplication and
$\Delta$, so $B_{s-1}(f,f)\in\cC$. This quadratic form is
$\langle My,y\rangle$ for a positive semidefinite $M$ whose kernel
is exactly the cylindrical directions of $f$: by polarization,
$B_{s-1}(f,f)(v)=0$ if and only if $\partial_vf\equiv0$.
Essentiality of $f$ therefore makes $M$ positive definite. Since
\[
 \Gamma(\langle My,y\rangle,\langle M^jy,y\rangle)
       =4\langle M^{j+1}y,y\rangle,
\]
all positive powers of $M$ give quadratic forms in $\cC$.
Choose a polynomial $p$ with $p(0)=0$ and $p(\lambda)=1$ at every
eigenvalue of $M$. Then $p(M)=\Id$ gives $q\in\cC$.
Consequently $\cA\subseteq\cC$, proving
\eqref{eq:smallest-clifford-pullback}.

At a point where $Q<R^2$, the Gram matrix of
$\nabla R,\nabla P_1,\ldots,\nabla P_m$ has determinant \(4^{m+1}R^{m-1}(R^2-Q)>0\).
Thus $(R,P)$ has an open set in its
image, and $R$ is algebraically independent of the pulled-back
polynomial ring. This proves the asserted isomorphism. For a finitely
generated connected graded algebra, the minimal homogeneous generator
number is $\dim(\cA_+/\cA_+^2)$, where $\cA_+$ is its positive-degree
ideal. Adjoining $t$ adds exactly one independent class.
\end{proof}

When $Q=R^2$ identically, the relation $q(P)=R^2$ must be retained;
the conclusion about one additional generator does not apply.

Choose minimal homogeneous generators
$\cA_q(f)=\R[u_1,\ldots,u_r]$, with $\deg u_i=d_i$, and write
\[
 \Delta u_i=b_i(u),\qquad \Gamma(u_i,u_j)=c_{ij}(u),\qquad q=q_*(u).
\]
The lifted variables $U_i=u_i(P)$ have the differential table
\begin{equation}\label{eq:clifford-minimal-table}
\begin{gathered}
 \Delta R=2N,\qquad \Gamma(R,R)=4R,\qquad
 \Delta U_i=4R\,b_i(U),\\
 \Gamma(R,U_i)=4d_iU_i,\qquad
 \Gamma(U_i,U_j)=4R\,c_{ij}(U).
\end{gathered}
\end{equation}
Thus the target table determines the calculation for any potential
$G=R^a\phi(U)$; only $\Delta R$ depends on the source dimension.
For a target potential factoring through $\cA_q(f)$ as $H=\phi(u)$,
this applies when $\phi$ is $C^3$ on an open neighborhood of
$u(\R^m\setminus Z)$.

\begin{proposition}[Target equivalence and dimension monotonicity]
\label{prop:clifford-transfer-monotonicity}
Fix a Clifford map $P_0:\R^{N_0}\to\R^m$ and, for $k\ge2$, set
\[
 P_k(x_1,\ldots,x_k)=\sum_{\alpha=1}^kP_0(x_\alpha),
 \qquad N_k=kN_0,\qquad R=\sum_{\alpha=1}^k|x_\alpha|^2.
\]
Assume the target hypotheses on $f,Z,H,M$ of
\cref{thm:clifford-area-transfer}, except for $H\cL(H)\ge0$.
For fixed $a>0$, put $G_k=R^aH(P_k)$. Then
\[
 H\cL(H)\ge0\text{ off }Z
 \quad\Longleftrightarrow\quad
 G_k\cL(G_k)\ge0\text{ off }P_k^{-1}(Z)
 \text{ for some }k\ge2.
\]
Once the source inequality holds for one such $k$, it holds for every
larger $k$. In particular, the target inequality is equivalent to the
source inequalities in all sufficiently large members of this family.
\end{proposition}

\begin{proof}
The image of $(R,P_k)$ is exactly \(\{(z,y)\in\R\times\R^m:z\ge|y|\}\).
Indeed, $|P_k|\le R$. Conversely, choose a unit vector $e$ with
$y=|y|e$ (arbitrary if $y=0$), and write
$y=y_1+y_2$, where
$y_1=(z+|y|)e/2$ and $y_2=-(z-|y|)e/2$.
The eigenvector construction in
\cref{lem:Clifford-algebra-pullback} realizes $P_0(x_i)=y_i$
with $|x_i|^2=|y_i|$. Two blocks therefore suffice.

For $y\notin Z$ and $z\ge|y|$, put $b=a(a+2\ell)$ and
\[
 \Psi(z,y)=z^aH(y),\qquad
 \mathcal W(z,y)=4z^{2a+1}W_H(y)+4bz^{2a-1}H(y)^2.
\]
For $N=N_k$, let $\mathcal D_N(z,y)$ be the value of
$G_k\cL(G_k)$ at any point with $(R,P_k)=(z,y)$.
Formula \eqref{eq:clifford-transfer-defect} shows that this value
is well-defined, and \eqref{eq:clifford-transfer-gradient} gives
$|\nabla G_k|^2=\mathcal W(z,y)>0$ there.
For two block-sum dimensions $N'=N_{k'}\ge N=N_k$, the same
formula gives the finite difference
\begin{equation}\label{eq:clifford-dimension-monotonicity}
 \mathcal D_{N'}(z,y)-\mathcal D_N(z,y)
 =\frac{2a(N'-N)}{z}\Psi(z,y)^2\mathcal W(z,y)\ge0.
\end{equation}
The common image therefore proves monotonicity of the source
inequalities.

The target inequality implies the source inequality for large $k$
by \cref{thm:clifford-area-transfer}. Conversely, fix $y\notin Z$.
The image calculation permits $(R,P_k)=(z,y)$ for every $z\ge|y|$.
By \eqref{eq:clifford-transfer-defect},
\[
 \lim_{z\to\infty}
 \frac{\mathcal D_{N_k}(z,y)}{16z^{4a+2}}
     =H(y)\cL(H)(y).
\]
If the source inequality holds, each expression on the left is
nonnegative. This proves necessity.
\end{proof}

Composition preserves umbilic maps
\cite[Proposition~5.1.8(1)]{RiedlerThesis25}. Induction from
\cref{thm:clifford-area-transfer,thm:smallest-clifford-pullback}
gives the following conclusions.

\begin{corollary}[Iteration]
\label{cor:iterated-clifford-transfer}
Let $f$ be a nonzero essential homogeneous polynomial of degree
$s\ge2$ on the target $\R^m$, and suppose $f,Z,H,M$ satisfy
\cref{thm:clifford-area-transfer}, with $H=\phi(u)$ as above.
Successive Clifford block sums with at least two blocks and
sufficiently large source dimensions yield minimizing phase boundaries.
After $j$ steps the defining degree is $2^js$ and the smallest algebra
is isomorphic to $\cA_q(f)[t_1,\ldots,t_j]$, with minimal generator
degrees
\[
 2^jd_1,\ldots,2^jd_r,\qquad 2,4,\ldots,2^j.
\]
Thus its minimal homogeneous generator number is exactly $r+j$.
The same ratio constant $M$ can be used at every step. If both signs
of $f$ occur near every zero, the minimizing supports are the complete
successive zero sets.
\end{corollary}

\begin{proof}
The ratio estimate \eqref{eq:clifford-transfer-ratio} preserves $M$.
The hypotheses on regularity, signs, the exceptional set, and the
nonvanishing of the gradient persist by the proof of \cref{thm:clifford-area-transfer};
the new potential remains a function of the new generators.
If $\partial_v(g\circ P)\equiv0$ for a nonzero homogeneous $g$
of degree $d>0$, then $v=0$: otherwise
$(g\circ P)(v+x)=(g\circ P)(x)$ gives vanishing order $2d$ at $v$,
whereas the submersion property of $P$ makes this equal to the
vanishing order of $g$ at $P(v)$, which is at most $d$.
The common-image calculation above gives $Q<R^2$ somewhere at each
step, so the algebra and degree assertions follow inductively from
\cref{thm:smallest-clifford-pullback}. At each stage a symmetric
Clifford system exists on $\Lambda^*\R^m$, with generators given by
exterior multiplication plus contraction; direct sums make its source
dimension arbitrarily large. Thus each required dimension bound can
be met. The support assertion follows at every stage from
\cref{thm:clifford-area-transfer}.
\end{proof}

